%
% merohomolog.tex -- Homologie für eine endliche Menge von Punkten
%
% (c) 2021 Prof Dr Andreas Müller, OST Ostschweizer Fachhochschule
%
\documentclass[tikz]{standalone}
\usepackage{amsmath}
\usepackage{times}
\usepackage{txfonts}
\usepackage{pgfplots}
\usepackage{csvsimple}
\usetikzlibrary{arrows,intersections,math,calc}
\definecolor{darkred}{rgb}{0.8,0,0}
\begin{document}
\def\skala{1}
\def\punkt#1{
	\fill[color=white] #1 circle[radius=0.06];
	\draw #1 circle[radius=0.06];
}
\def\ohr#1#2#3#4{
	\draw[color=blue,line width=1.2pt] #1 circle[radius=0.45cm];
	\draw[color=blue,line width=3.6pt,shorten <= 0.45cm] #1 -- #2;
	\draw[color=gray!20,line width=1.2pt] #1 -- #2;
	\node at #1 [above] {$z_{#3}$};
	\node[color=blue] at ($#1+(#4:0.65)$) {$\gamma_{#3}$};
	\punkt{#1}
}
\begin{tikzpicture}[>=latex,thick,scale=\skala]

\fill[color=gray!20,rounded corners=0.8cm] (-6,-3.5) rectangle (6,3.5);
\node at (-5.5,3) {$U$};

\coordinate (A) at (4,1);
\coordinate (A1) at (5,1);

\coordinate (B) at (3,-2);
\coordinate (B1) at (3,-3);

\coordinate (C) at (1,0.5);
\coordinate (C1) at (1.5,3);

\coordinate (D) at (0.5,-1);
\coordinate (D1) at (1.0,-3);

\coordinate (E) at (-0.5,2);
\coordinate (E1) at (-1.0,3);

\coordinate (F) at (-2,0.5);
\coordinate (F1) at (-5,1);

\coordinate (G) at (-2.5,-0.5);
\coordinate (G1) at (-3.0,-3);

\begin{scope}
	\clip[rounded corners=0.8cm] (-5,-3) rectangle (5,3);

	\ohr{(A)}{(A1)}{7}{110}
	\ohr{(B)}{(B1)}{6}{130}
	\ohr{(C)}{(C1)}{5}{240}
	\ohr{(D)}{(D1)}{4}{120}
	\ohr{(E)}{(E1)}{3}{-70}
	\ohr{(F)}{(F1)}{2}{-20}
	\ohr{(G)}{(G1)}{1}{170}
\end{scope}

\draw[color=darkred,rounded corners=1.2cm] (-5,-3) rectangle (5,3);

\node[color=darkred] at (0,3) [above] {$\gamma$};


\draw[<-,color=darkred,line width=1.2pt] (0,-3) -- ++(-0.1,0);
\draw[<-,color=darkred,line width=1.2pt] (-5,0) -- ++(0,0.1);
\draw[<-,color=darkred,line width=1.2pt] (5,0) -- ++(0,-0.1);


\end{tikzpicture}
\end{document}

